\documentclass[10pt]{article}
\usepackage[T1]{fontenc}
\usepackage[utf8]{inputenc}
\usepackage[french]{babel}
\usepackage{amsmath,amssymb,mathrsfs,mathpazo}
\usepackage{geometry}
\geometry{a4paper,margin=2cm}
\usepackage{graphicx,array,tabularx,tabularray,colortbl,multirow,multicol,booktabs,enumitem,pifont,ulem,fancybox,diagbox,cancel}
\usepackage{pstricks,pstricks-add,pst-plot,pst-tree,pst-node,pst-text,pst-func,pst-math,pst-bspline,pst-3dplot,pst-eucl,pst-solides3d}
\let\pscalculate\undefined
\usepackage{pst-bezier}
\usepackage{tikz,pgfplots,tkz-tab}
\usetikzlibrary{arrows,arrows.meta,patterns,calc,positioning,trees,decorations.markings,shapes.geometric,shapes.misc,shapes.arrows,intersections,decorations.pathmorphing,decorations.pathreplacing,backgrounds,shadows,babel}
\pgfplotsset{compat=1.18}
\usepgfplotslibrary{fillbetween}
\usepackage[np]{numprint}
\usepackage{listings,makecell,pas-tableur,twoopt}
\let\barre\undefined
\usepackage[pstricks]{tablvar}
\definecolor{darkgreen}{RGB}{0,100,0}
\providecommand{\red}{\color{red}}
\providecommand{\blue}{\color{blue}}
\providecommand{\green}{\color{green}}
\providecommand{\black}{\color{black}}
\providecommand{\white}{\color{white}}
\providecommand{\gray}{\color{gray}}
\providecommand{\cyan}{\color{cyan}}
\providecommand{\magenta}{\color{magenta}}
\providecommand{\yellow}{\color{yellow}}
\providecommand{\violet}{\color{violet}}
\providecommand{\orange}{\color{orange}}
\providecommand{\euro}{\texteuro}
\providecommand{\EUR}{\texteuro}
\providecommand{\ieme}{\textsuperscript{e}}
\providecommand{\degres}{\ensuremath{^\circ}}
\providecommand{\degree}{\ensuremath{^\circ}}
\providecommand{\vv}[1]{\overrightarrow{#1}}
\providecommand{\Coord}[1]{#1}
\providecommand{\decosix}{}
\providecommand{\starredbullet}{\ensuremath{\bullet}}
\pagestyle{empty}
\setlength{\parindent}{0pt}
\usepackage{adjustbox}\DeclareUnicodeCharacter{2236}{:}\DeclareUnicodeCharacter{2190}{\ensuremath{\leftarrow}}
\providecommand{\R}{\mathbb{R}}
\providecommand{\N}{\mathbb{N}}
\providecommand{\Z}{\mathbb{Z}}
\providecommand{\Q}{\mathbb{Q}}
\providecommand{\D}{\mathbb{D}}
\providecommand{\C}{\mathbb{C}}
\providecommand{\e}{{\mathrm{e}}}
\providecommand{\vect}[1]{\overrightarrow{#1}}
\providecommand{\vectt}[1]{\overrightarrow{\mathrm{#1}}}
\providecommand{\vvt}[1]{\overrightarrow{\mathrm{#1}}}
\providecommand{\np}[2][]{{#2~#1}}
\providecommand{\vv}[1]{{\overrightarrow{#1}}}
\renewcommand{\d}{{\mathrm{d}}}
\providecommand{\degres}{{^\circ}}
\providecommand{\degre}{{^\circ}}
\providecommand{\blue}{\color{blue}}
\providecommand{\red}{\color{red}}
\providecommand{\enskip}{\quad}
\providecommand{\llbracket}{\lbrack\!\lbrack}
\providecommand{\rrbracket}{\rbrack\!\rbrack}
\providecommand{\euro}{{\text{€}}}
\providecommand{\no}{{\text{n°}}}
\providecommand{\interval}[2]{{[#1;#2]}}
\providecommand{\intervFO}[2]{{[#1;#2[}}
\providecommand{\dotfill}{\ldots}

\providecommand{\ds}{\displaystyle}
\providecommand{\textding}[1]{#1}
\providecommand{\Oij}{\(\mathrm{O};\vec{\imath},\vec{\jmath}\)}
\providecommand{\Oijk}{\(\mathrm{O};\vec{\imath},\vec{\jmath},\vec{k}\)}
\providecommand{\Ouv}{\(\mathrm{O};\vec{u},\vec{v}\)}
\providecommand{\intervFF}[2]{\ensuremath{[#1;#2]}}
\renewcommand{\euro}{\ifmmode\text{\texteuro}\else\texteuro\fi}
\DeclareUnicodeCharacter{00B0}{\ifmmode^{\circ}\else\textdegree\fi}
\DeclareUnicodeCharacter{20AC}{\ifmmode\text{\texteuro}\else\texteuro\fi}
% Énoncés : Éducation nationale. Transcription et corrigés : APMEP, Denis Vergès et contributeurs. Découpage et classement : CoopMaths, Sébastien Lozano, Éric Elter et contributeurs. Adaptation de la présentation : Maths Explorer.
% CC BY-SA : https://coopmaths.fr/www/about/
% https://coopmaths.fr/alea/static/bac/2022/tex/bac_2022_09_sujet1_metropole_3.tex
\begin{document}
\noindent\textit{APMEP / CoopMaths. CC BY-SA. Présentation adaptée par Maths Explorer.}\par\medskip
\bigskip

\textbf{Partie A}

\medskip

On considère la fonction $f$ définie sur l'intervalle $[1~;~+\infty[$ par 

\[f(x) = \dfrac{\ln x}{x},\]

où $\ln$ désigne la fonction logarithme népérien.



\begin{enumerate}
\item Donner la limite de la fonction $f$ en $+\infty$.
\item On admet que la fonction $f$ est dérivable sur l'intervalle $[1~;~+\infty[$ et on note $f'$ sa fonction dérivée.
	\begin{enumerate}
		\item Montrer que, pour tout nombre réel $x \geqslant 1,\:f'(x) = \frac{1 - \ln x}{x^2}$.
		\item Justifier le tableau de signes suivant, donnant le signe de $f'(x)$ suivant les valeurs de $x$.
\begin{center}
\psset{unit=1cm}
\begin{pspicture}(6,1)
\psframe(6,1)
\psline(0,0.5)(6,0.5)\psline(2,0)(2,1)
\uput[u](1,0.4){$x$}\uput[u](2.1,0.4){$1$}\uput[u](4,0.4){e}\uput[u](5.5,0.4){$+ \infty$}
\rput(1,0.25){$f'(x)$}\rput(3,0.25){$+$}\rput(4,0.25){0} \rput(5,0.25){$-$}
\end{pspicture}
\end{center}
		\item Dresser le tableau de variations complet de la fonction $f$.
	\end{enumerate}
\item Soit $k$ un nombre réel positif ou nul.
	\begin{enumerate}
		\item Montrer que, si $0 \leqslant k \leqslant \dfrac{1}{\text{e}}$, l'équation $f(x) = k$ admet une unique solution sur l'intervalle $[1~;~\text{e}]$.
		\item Si $k > \dfrac{1}{\text{e}}$, l'équation $f(x) = k$ admet-elle des solutions sur l'intervalle $[1~;~+\infty[$ ?

Justifier.
	\end{enumerate}
\end{enumerate}

\bigskip

\textbf{Partie B}

\medskip

Soit $g$ la fonction définie sur $\R$ par : 

\[g(x) = \text{e}^{\frac{x}{4}}.\]

On considère la suite $\left(u_n\right)$ définie par $u_0 = 1$ et, pour tout entier naturel $n$ :

$u_{n+1} = \text{e}^{\frac{u_n}{4}}$ c'est-à-dire : $u_{n+1} = g\left(u_n\right)$.

\medskip

\begin{enumerate}
\item Justifier que la fonction $g$ est croissante sur $\R$.
\item Montrer par récurrence que, pour tout entier naturel $n$, on a : $u_n \leqslant  u_{n+1} \leqslant \text{e}$.
\item En déduire que la suite $\left(u_n\right)$ est convergente.
\end{enumerate}

On note $\ell$ la limite de la suite $\left(u_n\right)$ et on admet que $f$ est solution de l'équation : 
\[\text{e}^{\frac{x}{4}} = x.\]

\begin{enumerate}[resume]
\item En déduire que $\ell$ est solution de l'équation $f(x) = \dfrac14$, où $f$ est la fonction étudiée dans
la partie A.
\item Donner une valeur approchée à $10^{-2}$ près de la limite $\ell$ de la suite $\left(u_n\right)$.
\end{enumerate}

\bigskip

\end{document}
